%% LyX 2.4.0~RC3 created this file.  For more info, see https://www.lyx.org/.
%% Do not edit unless you really know what you are doing.
\documentclass[12pt,british]{article}
\usepackage{mathptmx}
\renewcommand{\familydefault}{\rmdefault}
\usepackage[LGR,T1]{fontenc}
\usepackage{textcomp}
\usepackage[latin9]{inputenc}
\usepackage[a4paper]{geometry}
\geometry{verbose,tmargin=2cm,bmargin=2cm,lmargin=2cm,rmargin=2cm}
\setcounter{secnumdepth}{5}
\setcounter{tocdepth}{5}
\usepackage{graphicx}
\usepackage[authoryear]{natbib}

\makeatletter

%%%%%%%%%%%%%%%%%%%%%%%%%%%%%% LyX specific LaTeX commands.
\DeclareRobustCommand{\greektext}{%
  \fontencoding{LGR}\selectfont\def\encodingdefault{LGR}}
\DeclareRobustCommand{\textgreek}[1]{\leavevmode{\greektext #1}}


%%%%%%%%%%%%%%%%%%%%%%%%%%%%%% User specified LaTeX commands.
\usepackage{bibentry}
\nobibliography{numerics}

\makeatother

\usepackage{babel}
\begin{document}
\begin{itemize}
\item \bibentry{ECMWF_aifs23}

``Do we want to leverage our knowledge of spectral transforms and
build something like NVIDIA's Neural Operator system, FourCastNet?
Do we want to leverage the vast research that goes into vision transformers
and build something like Pangu-Weather, FengWu or FuXi? Or do we want
to utilise the grid-flexibility and parameter efficiency of Graph
Neural Networks like the work of Ryan Keisler or Google Deepmind's
GraphCast?''

Chosen ``Graph Neural Networks''
\begin{itemize}
\item using reduced Gaussian grids
\item ``13 pressure levels, runs at approximately 1 degree resolution and
makes predictions for wind, temperature, humidity and geopotential.
At the surface AIFS makes predictions for 2 m temperature, 10 m winds,
surface pressure and more.''
\item ``trained to minimise mean squared error''
\end{itemize}
\item \bibentry{ECMWF_aifsEns25}
\begin{itemize}
\item AIFS Single released in February.
\item ``The new ensemble model outperforms state-of-the-art physics-based
models for many measures, including surface temperature, with gains
of up to 20\%. At the moment, it works at a lower resolution (31 km)
than the physics-based ensemble system (9 km), which remains indispensable
for high-resolution fields and coupled Earth-system processes.''
\item ECMWF is therefore also exploring hybrid systems that leverage the
strengths of both approaches.
\item Anemoi framework 
\end{itemize}
\item \bibentry{ECMWF_aifs25}
\begin{itemize}
\item ``AIFS outperforms state-of-the-art physics-based models for many
measures, including tropical cyclone tracks, with gains of up to 20\%.''
\item ``increased speed and a reduction of approximately 1,000 times in
energy use for making a forecast''.
\end{itemize}
\item \bibentry{LAC+24}
\begin{itemize}
\item ``graph neural network (GNN) encoder and decoder, and a sliding window
transformer processor''
\item ``AIFS is implemented in Python and utilises the PyTorch (Paszke
et al. {[}2019{]}) and PyTorch Geometric (Fey and Lenssen {[}2019{]})
machine learning frameworks. Together with AIFS we develop the Anemoi
(Greek: \textquoteright Winds\textquoteright ) framework, a toolbox
that aims to provide building blocks for data driven weather forecast
models.''
\item ``A diverse set of deep learning architectures is being employed
in the context of data-driven weather forecasts, e.g., vision transformers
(Bi et al. {[}2023{]}, Nguyen et al. {[}2023{]}), neural operators
(Pathak et al. {[}2022{]}) and graph neural networks (GNNs, Keisler
{[}2022{]}, Lam et al. {[}2023{]}).''
\item ``The first version of AIFS was introduced in experimental operational
mode in October 2023 (Lang et al. {[}2024{]}). Its architecture was
inspired by Keisler {[}2022{]} and Lam et al. {[}2023{]}. Specifically,
it consisted of an encoder-processor-decoder design similar to that
described in Battaglia et al. {[}2016{]}.''
\item ``The current implementation of AIFS shares certain limitations with
some of the other data-driven weather forecast models that are trained
with a weighted MSE loss, such as blurring of the forecast fields
at longer lead times. The blurring can be traced back to the MSE training
objective, where the model learns an implicit smoothing in order to
avoid the double penalty (e.g. Hoffman et al. {[}1995{]}, Ebert et
al. {[}2013{]}) of misplaced structures in the forecast.''
\item ``The strength and weaknesses of data-driven forecast models have
been investigated in Ben Bouall�gue et al. {[}2024{]} and Charlton-Perez
et al. {[}2024{]}, and in an idealised context in Hakim and Masanam
{[}2023{]}''
\end{itemize}
\item \bibentry{CDD+24}
\begin{itemize}
\item ``The four machine learning models considered (FourCastNet, Pangu-Weather,
GraphCast and FourCastNet-v2) produce forecasts that accurately capture
the synoptic-scale structure of the cyclone including the position
of the cloud head, shape of the warm sector and location of the warm
conveyor belt jet, and the large-scale dynamical drivers important
for the rapid storm development such as the position of the storm
relative to the upper-level jet exit.''
\item ``These models have all been shown to produce skillful 0-10 day forecasts
of the 500 hPa geopotential height field, based on the widely used
Anomaly Correlation Coefficient metric. All four models use an encode-process-decode
framework but with differing architectures:
\begin{itemize}
\item FourCastNet, developed by NVIDIA and based on Fourier Neural Operators
(FNO) with a vision transformer architecture;
\item FourCastNet version 2, which builds on FourCastNet by using spherical
FNOs;
\item Pangu-Weather, developed by Huawei and based on a three-dimensional
Earth-specific transformer and hierarchical temporal aggregation (four
deep neural networks with different lead times)
\item GraphCast, developed by Google DeepMind and based on graph neural
networks.
\end{itemize}
\item ``The ML models are all autoregressive, which means model output
from a given time step can be used to predict output at the next time
step.''
\item Good summary of the methods on page 8.
\end{itemize}
\item Cyril Morcrette -- Machine learning for physical processes
\begin{itemize}
\item CARAMEL -- cloud resolving model machine learning
\item ENNUF -- easy neural networks for the UM in Fortran
\item for a warmer world, train based on RH rather than on $q_{v}$
\end{itemize}
\item \bibentry{DLL+24}
\begin{itemize}
\item ``The most widely used ML models in rainfall forecasting include
support vector machine (SVM), random forest (RF), artificial neural
networks (ANN), and their variants.''\\
This is about models that predict just rainfall, not the state of
the whole atmosphere.
\end{itemize}
\item \bibentry{KSO+25}
\begin{itemize}
\item ML used to correct WRF output.
\end{itemize}
\item \bibentry{KFB13}
\begin{itemize}
\item ``The newly developed NN convection parameterization has been tested
in National Center of Atmospheric Research (NCAR) Community Atmospheric
Model (CAM). It produced reasonable and promising decadal climate
simulations for a large tropical Pacific region.''
\end{itemize}
\item \bibentry{dBL23}
\item \bibentry{KSA+21}
\begin{itemize}
\item ``we use end-to-end deep learning to improve approximations inside
computational fluid dynamics for modeling two-dimensional turbulent
flows. For both direct numerical simulation of turbulence and large-eddy
simulation, our results are as accurate as baseline solvers with 8
to 10$\times$ finer resolution''
\item ``stable during long simulations and generalizes to forcing functions
and Reynolds numbers outside of the flows where it is trained''
\item ``uses discrete equations that give pointwise accurate solutions
on an unresolved grid.''
\item ``The model learns how to interpolate local features of solutions''\\
\includegraphics[width=1\textwidth]{figures/KSA+21_fig1}
\item ``We train the model inside a standard numerical method for solving
the underlying PDEs as a differentiable program, with the neural networks
and the numerical method written in a framework {[}JAX (38){]} supporting
reverse-mode automatic differentiation.''.
\item ``we wrote a new CFD code in JAX''.
\item ``the convective flux model improves the approximation of the discretized
convection operator, the divergence operator enforces local conservation
of momentum according to a finite volume method, and the pressure
projection enforces incompressibility''.
\item ``two types of ML components: learned interpolation and learned correction''.\\
``The equation maintains the same symmetries and scaling properties
(e.g., rescaling coordinates x \textrightarrow{} \textgreek{\textlambda}x)
as the original equations and 2) the interpolation is always at least
first-order accurate with respect to the grid spacing, by constraining
the filters to sum to unity''.
\end{itemize}
\item \bibentry{KYL+24}
\begin{itemize}
\item ``NeuralGCM can accurately track climate metrics for multiple decades,
and climate forecasts with 140-kilometre resolution show emergent
phenomena such as realistic frequency and trajectories of tropical
cyclones''.\\
``Unlike physical models, machine-learning models misrepresent derived
(diagnostic) variables such as geostrophic wind''\\
\includegraphics[width=1\textwidth]{figures/KYL+24_fig1}
\item ``NeuralGCM, a fully differentiable hybrid GCM of Earth\textquoteright s
atmosphere. NeuralGCM is trained on forecasting up to 5-day weather
trajectories sampled from ERA5. Differentiability enables end-to-end
\textquoteleft online training'{}''.
\item ``NeuralGCM demonstrates that incorporating machine learning is a
viable alternative to building increasingly detailed physical models32
for improving GCMs.''
\end{itemize}
\item \bibentry{QHZ+25}\ ``FourCastNet ... shows promise in learning from
strong storms in one region and forecasting them in another region''
\begin{itemize}
\item FourCastNet (1), Pangu (2), GraphCast (3), AIFS (4), Fuxi (5), and
Stormer (6) ... ``deep neural networks (NNs) that predict the evolution
of the 3D global atmospheric state in six or twelve hour increments
after being trained on the ERA5 reanalysis dataset from 1979 to 2015''
\item ``Whether trained with the Full or noTC dataset, FourCastNet is unable
to reproduce gradient-wind balance, a key physical constraint that
is obeyed by TCs in the training set (Fig. 5). Enforcing gradient-wind
balance in the loss function could potentially help the model learn
this physical relationship, though there could be tradeoffs in the
accuracy or the representation of other physical properties (e.g.,
the ageostrophic flow).''
\end{itemize}
\item \bibentry{AJM+25}\ \\
PDEs discontinuous solutions
\begin{itemize}
\item PINNS -- Physics-Informed Neural Networks (PINNs)
\end{itemize}
\begin{enumerate}
\item Physics-modification (PM) methods improve accuracy by modifying the
system\textquoteright s physics, such as adding artificial viscosity
or enforcing entropy constraints.
\item Loss and training modification (LM) techniques focus on regularizing
the loss landscape, often by refining the loss term in high-error
regions. 
\item Architecture-modification (AM) propose advanced network designs to
handle discontinuities better. 
\begin{itemize}
\item the generalization error in the solutions quickly increased
\end{itemize}
\end{enumerate}
\begin{itemize}
\item ``PINNs integrate physical principles into the neural network's training
process, by adding the differential equations to a composite loss
function''. The training error is the difference between the NN solution
and the exact solution. The loss function is the training error.
\end{itemize}
\item \bibentry{AMK+25}\ \\
This an application of PINNS. See RPK19
\item \bibentry{RPK19}\ \\
data-driven solution and data-driven discovery of PDEs
\begin{itemize}
\item continuous time
\begin{itemize}
\item new family of data-efficient spatio-temporal function approximators
\end{itemize}
\item discrete time models. 
\begin{itemize}
\item implicit RK time stepping
\end{itemize}
\end{itemize}
\item \bibentry{MMN20}
\item \bibentry{YO20}
\item \bibentry{BRPG19}
\item \bibentry{LFZ+24}
\item \bibentry{SSCG}\ \\
``FuXi (Chen et al., 2023), a state-of-the-art AIWP model, is selected
as an example to implement conservation schemes. The base architecture
of FuXi is Swin-Transformer, which has been adopted by a wide range
of AIWP models (e.g., Nguyen, Shah, et al., 2023; Willard et al.,
2024), making FuXi an ideal candidate for showcasing the integration
of physical constraints\\
Solves ML mode and conservation laws. ``The specific total water,
total precipitation, and air temperature constraints are applied using
multiplicative ratios across all grid cells. For specific total water,
constraints are limited to pressure levels below 600 hPa, as the moisture
content in the upper atmosphere is significantly lower, and enforcing
constraints at higher levels would lead to strong outlier effects
within the z-score space.''
\end{itemize}
\bibliographystyle{abbrvnat}
\bibliography{/home/hilary/latex/lit/ML}

\end{document}
